

\begin{table*}[t]
    \centering
    \footnotesize
    \setlength{\tabcolsep}{2pt}  % Reduce horizontal padding
    \begin{minipage}[t]{0.48\textwidth}
    \centering
    \begin{tabular}{|>{\centering\arraybackslash}m{2.3cm}|*{4}{c|}c|c|c|c|}
    \hline
    \multirow{2}{*}[-5.5ex]{\begin{tabular}[c]{@{}c@{}}\textbf{Enclave}\\\textbf{Platforms}\end{tabular}} & 
    \multicolumn{4}{c|}{\textbf{Security}} & \multicolumn{1}{c|}{\textbf{Usability}} & \multicolumn{3}{c|}{\textbf{Implementation}} \\
    \cline{2-9}
    & \rotatebox{90}{\parbox{1.5cm}{\centering SC LLC}} &
    \rotatebox{90}{\parbox{1.5cm}{\centering SC Others}} &
    \rotatebox{90}{\parbox{1.5cm}{\centering TES}} &
    \rotatebox{90}{\parbox{1.5cm}{\centering Exec-Time}} &
    \rotatebox{90}{\parbox{1.5cm}{\centering Shared Memory}} &
    \rotatebox{90}{\parbox{1.5cm}{\centering HW}} &
    \rotatebox{90}{\parbox{1.5cm}{\centering SW. Infra.}} &
    \rotatebox{90}{\parbox{1.5cm}{\centering Apps}} \\
    \hline
    Commercial \cite{mckeen2013sgx, costan2016intel,cheng2024intel, sev2020strengthening,trustzone04, pinto2019demystifying, li2022design}& \red{V} & \red{V} & \red{V} & \red{V} & \green{Yes} & \green{TO} & \green{Yes} & \green{Yes} \\
    \hline
    Academic \cite{azab2011sice, brasser2019sanctuary, deng2014equalvisor, feng2021scalable, van2023cheri, ferraiuolo2017komodo, weiser2019timber} & \red{V} & \red{V} & \red{V} & \red{V} & \green{Yes} & \orange{R/S} & \orange{Some} & \orange{Some} \\
    \hline
    Sanctum \cite{costan2016sanctum} & \green{P} & \red{V} & \red{V} & \red{V} & \green{Yes} & \red{Sim} & \red{No} & \red{No} \\
    \hline
    KeyStone \cite{lee2020keystone}, CURE\cite{bahmani2021cure}, Hybrid\cite{enclavecaches} & \green{P} & \red{V} & \red{V} & \red{V} & \green{Yes} & \green{RTL} & \green{Yes} & \green{Yes} \\
    \hline
    MI6 \cite{bourgeat2019mi6} & \green{P} & \green{P} & \green{P} & \red{V} & \red{No} & \green{RTL*} & \red{No} & \red{No} \\
    \noalign{\hrule height 1pt}
    \textbf{Citadel} & \textbf{\green{P}} & \textbf{\green{P}} & \textbf{\green{P}}  & \textbf{\red{V}} & \textbf{\green{Yes}} & \textbf{\green{RTL}} & \textbf{\green{Yes}} & \textbf{\green{Yes}} \\
    \noalign{\hrule height 1pt}
    \end{tabular}
    \caption{Comparison with Other Enclave Platforms.
    SC: Side Channel, TES: Transient Execution SC, Exec-Time: Execution-Time SC,
    P: Protected, V: Vulnerable, TO: Taped-out, Sim: Simulator, R/S: RTL or Sim, RTL*: Incomplete}
    \label{tab:enclave-platforms}
    \end{minipage}%
    \hfill
    \begin{minipage}[t]{0.48\textwidth}
    \centering
    \begin{tabular}{|>{\centering\arraybackslash}m{2.1cm}|*{4}{c|}*{2}{c|}c|c|}
    \hline
    \multirow{2}{*}[-5.5ex]{\begin{tabular}[c]{@{}c@{}}\textbf{Defense}\\\textbf{Mechanism}\end{tabular}} & 
    \multicolumn{4}{c|}{\textbf{Security}} & \multicolumn{2}{c|}{\textbf{Usability}} & \multicolumn{2}{c|}{\textbf{Overhead}} \\
    \cline{2-9}
    & \rotatebox{90}{\parbox{1.5cm}{\centering TES Cache}} &
    \rotatebox{90}{\parbox{1.5cm}{\centering TES Others}} &
    \rotatebox{90}{\parbox{1.5cm}{\centering Register Values}} &
    \rotatebox{90}{\parbox{1.5cm}{\centering TES Exec-Time}} &
    \rotatebox{90}{\parbox{1.5cm}{\centering Annotation Secrets}} &
    \rotatebox{90}{\parbox{1.5cm}{\centering Code Analysis}} &
    \rotatebox{90}{\parbox{1.5cm}{\centering Hardware}} &
    \rotatebox{90}{\parbox{1.5cm}{\centering Performance}} \\
    \hline
    ProSpeCT \cite{daniel2023prospect} & \green{P} & \green{P} & \green{P} & \green{P} & \red{Yes} & \red{CT} & \red{High} & 0-45\% \\
    \hline
    SPT \cite{choudhary2021speculative} & \green{P} & \green{P} & \green{P} & \green{P} & \green{No} & \red{CT} & \red{Sim} & 45\% \\
    \hline
    STT \cite{yu2019speculative,jauch2023secure} & \green{P} & \green{P} & \red{V} & \green{P} & \green{No} & \green{No} & \red{High} & 46\% \\
    \hline
    SpecShield \cite{barber2019specshield} & \green{P} & \green{P} & \red{V} & \green{P} & \green{No} & \green{No} & \red{Sim} & 21-73\% \\
    \hline
    NDA \cite{weisse2019nda} & \green{P} & \green{P} & \orange{S} & \green{P} & \green{No} & \green{No} & \red{Sim} & 100\% \\
    \hline
    SpectreGuard \cite{fustos2019spectreguard} & \green{P} & \green{P} & \green{P} & \green{P} & \red{Yes} & \green{No} & \red{Sim} & 8-20\% \\
    \hline
    ConTExT \cite{schwarz2020context} & \green{P} & \green{P} & \green{P} & \red{V} & \red{Yes} & \green{No} & \red{Sim} & 0-338\% \\
    \hline
    InvisiSpec \cite{yan2018invisispec} + \cite{khasawneh2019safespec, sakalis2019efficient, li2019conditional, ainsworth2020muontrap, saileshwar2019cleanupspec} & \green{P} & \red{V} & \green{P} & \red{V} & \green{No} & \green{No} & \red{Sim} & 21-72\% \\
    \hline
    DOLMA \cite{loughlin2021dolma} & \green{P} & \green{P} & \green{P} & \green{P} & \green{No} & \green{No} & \red{Sim} & 42.2\% \\
    \hline
    Sequential & \green{P} & \green{P} & \green{P} & \green{P} & \green{No} & \green{No} & \green{Negl} & 882\% \\
    \noalign{\hrule height 1pt}
    \textbf{Citadel} & \textbf{\green{P}} & \textbf{\green{P}} & \textbf{\green{P}} & \textbf{\red{V}} & \textbf{\green{No}} & \textbf{\orange{Min}} & \textbf{\green{Small}} & \textbf{\green{0\% - 5\%}} \\
    \noalign{\hrule height 1pt}
    \end{tabular}
    \caption{Comparison with Hardware Defense Mechanisms for Transient Execution Side Channels (TES), P: Protected, V: Vulnerable, S: Some, CT: Constant Time, Sim: Simulator}
     
    \label{tab:defense-mechanisms}
    \end{minipage}
    
\end{table*}

\section{Related Work}
\label{sec:relatedwork}
\noindent\textbf{Security of Enclave Platforms}
\label{sec:enclaveplatforms}
Table \ref{tab:enclave-platforms} illustrates the landscape of existing enclave platforms. 
All commercial platforms~\cite{mckeen2013sgx, costan2016intel,cheng2024intel, sev2020strengthening,trustzone04, pinto2019demystifying, li2022design} have proven vulnerable to TES attacks~\cite{chen2019sgxpectre,van2020sgaxe,gotzfried2017cache,brasser2017software,van2018foreshadow,van2020lvi,weisse2018foreshadow,ryan2019hardware,xu2024trustzonetunnel}.
Despite this, most commercial platforms and academic works on TEEs~\cite{azab2011sice, brasser2019sanctuary, deng2014equalvisor, feng2021scalable, van2023cheri, ferraiuolo2017komodo, weiser2019timber, van2023cheri, li2021confidential, lee2022cerberus, park2020nested, steinegger2021servas} do not include these attacks in their threat models, shifting the burden to developers and deferring defense to software analysis tools (see lower).
Although some academic proposals such as Sanctum \cite{costan2016sanctum}, KeyStone \cite{lee2020keystone}, and Cure \cite{bahmani2021cure} incorporate some defenses against cache side channels, only MI6 \cite{bourgeat2019mi6} formally addresses TES attacks but disables shared memory, failing to address \prob{} comprehensively.
Citadel builds on the open source code provided in MI6 \cite{bourgeat2019mi6} and reuses many of its microarchitectural isolation mechanisms.
The MI6 paper focused on evaluating the \emph{performance} of these mechanisms on an FPGA
by running SPEC benchmarks with no security guarantees by using Linux system calls as hooks to project the cost of a potential context switch to the SM.
Citadel's performance numbers for SPEC would be similar to MI6's evaluation.
Citadel stands out as the only platform that offers protection against TES attacks while maintaining shared memory usability and providing a complete implementation with hardware, software infrastructure, and applications.
\\

\noindent\textbf{Software Analysis Tools}
\label{sec:softwaremitigations}
Commercial TEEs provide speculation barriers and other flushing instructions for manual placement in software~~\cite{TrustDomainIntel, ArmSpecualtiveVulnerabilities, AMDsoftwarespecualtion}.
However, these countermeasures~\cite{mosier2024serberus, cauligi2022sok, disselkoen2020finding, guanciale2020inspectre} incur significant overhead, might not be sufficient \cite{barberis2022branch}, and their reliable automated application remains an open challenge~\cite{LLVMdevSpeculativeLoadHardening,IntelAnalysisSideChannel, MicrosoftMSVC,wang2019oo7}.

Other tools aim to target more precise speculative security policies~\cite{guarnieri2021hardware,cauligi2022sok}, but mostly target constant-time crypto~\cite{barthe2021high, cauligi2020constant, daniel2021hunting, guanciale2020inspectre, guarnieri2020spectector, patrignani2021exorcising, vassena2021automatically} or the sandboxing scenario~\cite{cheang2019formal, narayan2021swivel, qi2021spectaint, shen2018restricting, jenkins2020ghostbusting, oleksenko2020specfuzz, kirzner2021analysis} which makes them difficult to adapt to address \prob{} (see Section \ref{sec:beyondsandboxing}).
It is also unclear how pieces of code analyzed with these tools and satisfying \emph{stronger} speculative properties than RMI should be composed, or how they should communicate with each other.
In fact, introducing any universal read gadget in the code base would break security, even for the carefully analyzed speculative constant-time~\cite{cauligi2020constant} cryptographic library.
Citadel presents one strategy to compose different speculative policies.
It can even compose programs that guarantee the ``pre-Spectre'' non-interference properties~\cite{barthe2014system}, without worrying about speculation, as we can guarantee that predictors are independent of any secrets through isolation.
Finally, analysis tools struggle to scale to entire codebases, especially when trying to prove complex properties, further complicating this composition problem~\cite{guarnieri2020spectector,wang2019oo7,johannesmeyer2022kasper,wang2020kleespectre}.
Citadel addresses these limitations by isolating execution modes, enabling secure shared memory, and limiting software analysis to small code snippets.

Other tools like Revizor \cite{oleksenko2022revizor, garcia2023fuzzing} make it possible to fuzz CPUs for speculative leaks. 
To our knowledge, current available implementations only support X86. 
Adapting Revizor to RISC-V and RMI is left for future work.

\subsection{Hardware Mitigations for TES attacks}
\label{sec:exisitngTEAmech}
Numerous academic proposals \cite{ainsworth2020muontrap, cheang2019formal, canella2019systematic, barber2019specshield, fustos2019spectreguard, yan2018invisispec, yu2019speculative, saileshwar2019cleanupspec, weisse2019nda, khasawneh2019safespec, guarnieri2020spectector, guarnieri2021hardware, mosier2022axiomatic,ainsworth2021ghostminion, daniel2023prospect, choudhary2021speculative, kiriansky2018dawg, sakalis2019efficient, li2019conditional, narayan2023going} leverage hardware mechanisms to counter TES attacks.
Why not use some of these mechanisms as primitive and combine them with, say, MI6~\cite{bourgeat2019mi6} to address \prob{}?
As shown in Table \ref{tab:defense-mechanisms}, each mechanism offers a different trade-off between security, usability, and performance.
Within that design space, Citadel and its simple mechanisms can efficiently address \prob{}.

As explained in Section~\ref{sec:beyondsandboxing}, mechanisms that focus on the \emph{constant-time}~\cite{daniel2023prospect, choudhary2021speculative} and \emph{sandboxing}~\cite{yu2019speculative,jauch2023secure,barber2019specshield, yu2020speculative} scenarios are not the best building blocks to address \prob{} efficiently.
Similarly, NDA\cite{weisse2019nda} ignores Specter-BTB and does not prevent leakage of
register secrets through single transient micro-op.
Other Hardware Mechanisms such as Context\cite{schwarz2020context}, OISA\cite{yu2018data} and SpecterGuard\cite{fustos2019spectreguard} require extensive manual and error-prone secret annotations, on which \emph{security} depends.
In contrast, Citadel' code annotation is only required for \emph{performance} and any manual error would be caught by our static analysis tool.
Many other mechanisms could be a good fit to be used as primitives to build a platform such as Citadel.
Invisible speculation schemes such as InvisiSpec~\cite{yan2018invisispec} and others~\cite{ainsworth2020muontrap, khasawneh2019safespec, kiriansky2018dawg, li2019conditional, saileshwar2019cleanupspec, sakalis2019efficient, ainsworth2021ghostminion} only address cache side channels.
Similarly, DOLMA\cite{loughlin2021dolma} could be used to address \prob{} and even provide guarantees against the execution time SC (see Section~\ref{sec:execution_time_SC}), \emph{if} augmented with micro-architectural isolation similar to MI6.
\\

\noindent\textbf{Hardware and Performance Overhead}
However, even when considered as primitives, existing hardware mechanisms are \emph{complex} and incur significant \emph{overheads}.
Few of these mechanisms have been implemented in hardware (FPGAs or silicon), and many attacks that break these defenses \cite{yang2023pensieve} exploit details that are often overlooked in simulation, such as faithful memory hierarchy representation.
The two mechanisms that have been implemented incur significant hardware overhead.
STT on BOOM \cite{jauch2023secure} adds +15.7\% LUTs and 6.2\% FFs, while ProSpeCT \cite{daniel2023prospect} adds +17\% LUTs and +6\% FFs, an overhead potentially prohibitive for commercial deployment, especially given that they would require extra isolation mechanisms to address the TEE threat model (i.e., privilege attacker).
In that sense, these overheads should be compared with the minimal surface overhead of implementing Safe mode and Burst mode on top of MI6: +0.3\% LUTs and +1.0\% FFs.

These complex mechanisms also incur significant performance overhead.
DOLMA presents an overhead of +42.2\% on average, while NDA + 100\% and STT + 46\% for their most conservative threat models (but closest to our).
In contrast, Citadel shows a modest 0-5\% overhead on our benchmarks.